\ssscp{Distribution of terms for ‘opposite sex sibling’}{15-04-02-04}
\citet{GrimesCIP} shows the tight clustering of reflexes of
*bətaw ‘sister (male speaking)’
and *ñaRa ‘brother (female speaking)’ east of the Blust line,
and their scarcity west of the line,
as shown in \frf{15-04}.
Reflexes of *bətaw ‘sister (male speaking)’ are reported for
Ngaju Dayak and Berawan on Borneo meaning ‘sister’,
and for two Batak languages meaning ‘opposite sex sibling’.
\citet{Fox2011} notes that anthropologists
and native speakers have been unable to confirm the disputed
archaic reflex of \it{ɲahæ} reported for Ngaju Dayak,
but there may be trace evidence in oral epics in nearby languages that support it.
West of the Blust line reflexes of these two terms are quite
rare, with apparently none in the Philippines.
Yet the two etyma are assigned at the level of PMP,
based methodologically on an assumption of the validity of “WMP”
and “CMP” -- both of which are disputed or defunct subgroups (see
\frf{01-02}).

\begin{figure}[h]\centering
	\includegraphics[width=\columnwidth]{betaw-naRa.pdf}
	\fcp{Distribution of reflexes of *bətaw and *ñaRa}{15-04}
\end{figure}

\citet{BlustTrussel2020} acknowledge that *bətaw
and *ñaRa\footnote{
		Only the disputed Ngaju term shows the palatalised onset for *ñaRa.
		All languages in the east reflect *naRa,
		but this could be due to the common merger of *ñ/*n > \it{n}
		in the region.}
“are two of the key terms in determining the
history of marriage and social organization in PMP society.” These two sibling terms
figure prominently in the 35-year debate between Blust
and Fox regarding Malayo-Polynesian social structure.
One view of the data has two relative-age sibling terms for PMP:
*kaka ‘elder sibling’ and *huaji ‘younger sibling’.
These developed into four sibling terms around the beginning
of the contact zone (the Blust line) with slightly shifted meaning:
*kaka ‘elder sibling (same sex)’,
*huaji ‘younger sibling (same sex)’,
*bətaw ‘sister (male speaking)’
and *ñaRa ‘brother (female speaking)’.
The other view has all four terms assigned to PMP even though
there are no reflexes of the latter two in the Philippines
and much of western Indonesia.

East of the Blust line,
reflexes of both *bətaw
and *ñaRa are widespread, albeit uneven.
Languages investigated in the region that do not show cognates
with either term are represented by white circles in \frf{15-04}.

Some languages like Buru
and Lisela have reflexes of all four sibling terms: \it{kaka
{\tl} kai} ‘elder sibling (same sex)’,
\it{wai} ‘younger sibling (same sex)’,
\it{kaka-wai-t {\tl} kai-wai-t} ‘relatives’,
\it{feta} ‘sister (male speaking)’,
\it{naha, naha-t} ‘brother (female speaking)’
and a compound \it{feta-sar-naha} ‘reciprocally brother
and sister to each other’.
Some languages like Alune have two sibling terms: \it{wali} ‘same
sex sibling’ (from *huaji),
and \it{beta-i} ‘opposite sex sibling’ (from *bətaw).
Other languages like Luang have three sibling terms: \it{kakə-ni}
‘elder sibling’, \it{gar-ni} ‘younger sibling’, \it{nara} ‘opposite sex sibling’ (from *ñaRa).
See \citet{Fox2015} and \srf{07-02-03} for more variations
and the significance of these different patterns.

To recap, there is a concentration of reflexes of *bətaw
and *ñaRa in Wallacea.
The distribution of languages known to have both of these terms
is quite limited, often to lower-level branches within a subgroup.
Of our eight Wallacean subgroups,
only \tsc{Sula-Buru} has both these terms in all languages.
\tsc{Taliabu} also has all four,
even though no reflexes are known in other branches of \tsc{Celebic}.
Known reflexes of these two terms are extremely sparse in western
Indonesia, with only one language reported to have both terms.
There are no known reflexes in languages of the Philippines,
raising methodological questions about why these two terms are
reconstructed to the level of PMP.
This distributional pattern implies contact,
rather than inheritance from a common parent language.

There are many other words that have similar distribution patterns
that straddle the Blust line.
Some, like **mantəR ‘cuscus’ cannot be assigned to PMP.
Others like *bətaw and *ñaRa that have been assigned to PMP probably
should not have been,
due to their limited distribution
and the rethinking of PMP subgrouping.
Applying the same reasoning that was been used to propose PCEMP
on the basis of marsupial terms \citep{Blust1982,Blust2012b},
we should posit a subgroup that encompasses the range
of terms like **mantəR and **monsoŋ.
But because these words are found on both sides of the Blust
line, nobody is ready to do that,
since there is no accepted lower level subgroup based on the
historical phonology to which to assign them.
The uneven pockets of distribution of these terms not aligning
with other factors, mean that other considerations,
such as diffusion and contact, must be considered.

The distribution of individual words is not limited to subgroup boundaries.
This is not a new issue.
That is how words of Nahuatl origin (in Mexico) came through
Spanish becoming English, \it{coyote, tomato, chocolate} and \it{avocado}.
Lexical borrowings are a strong reason why subgroups should not
be based primarily on the lexicon.